\documentclass[12pt,a4paper]{report}

\usepackage{graphicx}
\usepackage{geometry}
\usepackage{setspace}
\usepackage{booktabs}
\usepackage{hyperref}
\usepackage{longtable}
\usepackage{float}
\usepackage{xcolor}
\usepackage{tikz}
\usepackage[utf8]{inputenc}
\usepackage[T1]{fontenc}

\geometry{margin=1in}
\setstretch{1.5}

% ── IMAGE PLACEHOLDER COMMAND ──────────────────────────────────────────────
% Usage: \imgbox{width_cm}{height_cm}{Caption label}
% To replace with real image: \includegraphics[width=0.8\textwidth]{filename.png}
\definecolor{PlaceholderBorder}{HTML}{0EA5E9}
\definecolor{PlaceholderFill}  {HTML}{EFF8FF}
\newcommand{\imgbox}[3]{%
  \begin{figure}[H]
    \centering
    \begin{tikzpicture}
      \draw[PlaceholderBorder, dashed, line width=1.2pt, rounded corners=6pt,
            fill=PlaceholderFill]
            (0,0) rectangle (#1,#2);
      \node[text=gray, font=\small\sffamily, align=center, text width=#1 cm]
            at (#1/2, #2/2) {\textbf{[Insert Screenshot / Diagram]}\\\medskip #3};
    \end{tikzpicture}
    \caption{#3}
  \end{figure}%
}

% ──────────────────────────────────────────────────────────────────────────

\begin{document}

%=============================================================
%  TITLE PAGE
%=============================================================
\begin{titlepage}
\centering

{\Large \textbf{B.P. Mandal College of Engineering, Madhepura}}\\
\vspace{0.3cm}
{\large (An Autonomous Institute, Affiliated to AKTU)}\\
\vspace{0.5cm}

Bachelor of Technology\\
in\\
Computer Science and Engineering

\vspace{1cm}

{\Huge \textbf{Project Report}}

\vspace{0.5cm}

{\LARGE \textbf{Decentralized E-Voting System}}\\[0.3cm]
{\large A Blockchain-Based Solution for Secure \& Transparent Elections}

\vspace{1cm}

Submitted By

\vspace{0.3cm}
\begin{center}
Ritesh Kumar (Roll No: XXXXXX)\\
Anuj Kumar   (Roll No: XXXXXX)\\
Ashish Kumar (Roll No: XXXXXX)\\
Shivam Kumar (Roll No: XXXXXX)\\
Deochandra Rishideo (Roll No: XXXXXX)
\end{center}

\vspace{1cm}

Under the guidance of

\vspace{0.2cm}
Prof. [Guide Name]\\
Assistant Professor\\
Department of Computer Science

\vfill

2022 -- 2026

\end{titlepage}

%=============================================================
%  DECLARATION
%=============================================================
\chapter*{Declaration}
\addcontentsline{toc}{chapter}{Declaration}

We hereby declare that the project entitled
\textbf{``Decentralized E-Voting System''}
is our own original work carried out at B.P. Mandal College of Engineering, Madhepura,
in partial fulfilment of the requirements for the award of the degree of
Bachelor of Technology in Computer Science and Engineering.

This work has not been submitted elsewhere for the award of any other degree or diploma.
All information and data presented in this report are true to the best of our knowledge.

\vspace{2cm}

\begin{tabular}{|c|l|l|l|}
\hline
S.No & Name               & Roll Number & Signature \\
\hline
1    & Ritesh Kumar       & XXXXXX      &           \\
\hline
2    & Anuj Kumar         & XXXXXX      &           \\
\hline
3    & Ashish Kumar       & XXXXXX      &           \\
\hline
4    & Shivam Kumar       & XXXXXX      &           \\
\hline
5    & Deochandra Rishideo & XXXXXX     &           \\
\hline
\end{tabular}

\vspace{2cm}

\noindent
Date: \underline{\hspace{4cm}} \\[0.4cm]
Place: \underline{\hspace{4cm}}

\newpage

%=============================================================
%  CERTIFICATE
%=============================================================
\chapter*{Certificate}
\addcontentsline{toc}{chapter}{Certificate}

This is to certify that the project report entitled

\begin{center}
\textbf{``Decentralized E-Voting System''}
\end{center}

is a bonafide project work carried out by the following students of the
7th Semester Bachelor of Technology (Computer Science and Engineering)
at B.P. Mandal College of Engineering, Madhepura, under my supervision
and guidance during the academic year 2025--2026.

\vspace{0.3cm}

\begin{center}
Ritesh Kumar, Anuj Kumar, Ashish Kumar, Shivam Kumar, Deochandra Rishideo
\end{center}

The project has been found satisfactory and is approved for submission.

\vspace{3cm}

\noindent
\begin{tabular}{p{6cm} p{6cm}}
Prof. [Guide Name]         & Prof. [HOD Name]\\
Project Guide              & Head of Department\\
Dept. of CSE               & Dept. of CSE\\
B.P. Mandal College of Engg. & B.P. Mandal College of Engg.\\
\end{tabular}

\vspace{2cm}

\noindent
Date: \underline{\hspace{4cm}} \\[0.4cm]
Place: Madhepura

\newpage

%=============================================================
%  ACKNOWLEDGEMENT
%=============================================================
\chapter*{Acknowledgement}
\addcontentsline{toc}{chapter}{Acknowledgement}

We take this opportunity to express our sincere gratitude to all those who
supported us throughout the development of this project.

First and foremost, we thank our project guide \textbf{Prof. [Guide Name]},
Assistant Professor, Department of Computer Science, for her invaluable guidance,
continuous encouragement, and constructive feedback at every stage of the project.

We extend our sincere thanks to \textbf{Prof. [HOD Name]}, Head of Department
of Computer Science, for providing the necessary resources and a conducive
environment for carrying out this work.

We are also grateful to the \textbf{Principal and Faculty Members} of
B.P. Mandal College of Engineering, Madhepura, for their administrative
support throughout the academic year.

Special thanks to the \textbf{Ethereum, Polygon, and Open-Source Blockchain communities}
whose publicly available tools, documentation, and frameworks made this project
possible.

Finally, we express our deepest gratitude to our \textbf{families and friends}
for their constant motivation, patience, and moral support during the entire journey.

\vspace{2cm}
\noindent
\hfill Ritesh Kumar, Anuj Kumar, Ashish Kumar, \\
\hfill Shivam Kumar, Deochandra Rishideo

\newpage

%=============================================================
%  ABSTRACT
%=============================================================
\chapter*{Abstract}
\addcontentsline{toc}{chapter}{Abstract}

The \textbf{Decentralized E-Voting System} is a blockchain-based web application
designed to provide a secure, transparent, and tamper-proof digital voting platform.
Traditional voting systems -- both physical and digital -- rely on centralized
databases that are vulnerable to manipulation, hacking, and administrative bias.

This project addresses these fundamental problems by leveraging the
\textbf{Ethereum-compatible Polygon blockchain} and \textbf{Solidity smart contracts}
to create a system where every vote is stored as an immutable on-chain transaction,
and the election logic is enforced by cryptographic code rather than a human authority.

The system automates the following critical operations:

\begin{itemize}
\item Election creation with configurable start and end times
\item Candidate registration by the election administrator
\item Voter whitelisting to prevent unauthorized participation
\item Ballot casting with strict one-vote-per-address enforcement
\item Automatic winner determination upon election close
\end{itemize}

The frontend is built using \textbf{Next.js}, connected to the blockchain via
\textbf{Ethers.js} and authenticated through \textbf{MetaMask} wallets.
Candidate and voter profile images are stored on \textbf{IPFS} using the
Pinata pinning service, ensuring decentralized media storage.

The smart contract was developed, tested, and deployed using the
\textbf{Hardhat} framework on the \textbf{Polygon Mumbai} testnet.

\newpage

\tableofcontents
\listoffigures
\listoftables

\newpage

%=============================================================
%  CHAPTER 1: INTRODUCTION
%=============================================================
\chapter{Introduction}

\section{Background of the Project}

Elections form the backbone of democratic societies. The integrity of any election
depends entirely on how trustworthy and transparent the voting mechanism is.
Conventional paper-based elections are slow, costly, and vulnerable to
booth-capturing or vote-tampering. Existing electronic voting systems, while faster,
still rely on centralized infrastructure that can be compromised.

Blockchain technology offers a revolutionary alternative. A blockchain is a
distributed, append-only ledger of records that no single party control.
Smart contracts -- self-executing programs deployed on a blockchain -- can enforce
election rules automatically and transparently without the need for a trusted intermediary.

This project presents a fully functional \textbf{Decentralized E-Voting DApp}
that implements these principles using widely adopted Web3 technologies.

\imgbox{13}{7}{Overall Project Overview Diagram -- show frontend, smart contract, and blockchain}

\section{Problem Statement}

Current digital voting systems suffer from the following limitations:

\begin{itemize}
\item \textbf{Centralized Single Point of Failure:} All votes are stored in a
      single database. If it is hacked or corrupted, the entire election is compromised.
\item \textbf{Lack of Voter Verifiability:} A voter cannot independently verify
      that their vote was recorded and counted correctly.
\item \textbf{Administrative Control:} The system administrator can potentially
      manipulate records before results are published.
\item \textbf{Opacity in Result Computation:} The process of tallying votes is
      hidden from the public, making it impossible to audit.
\item \textbf{Limited Accessibility:} Physical polling stations require voters
      to be physically present, limiting participation.
\end{itemize}

\section{Objectives}

The primary objectives of this project are:

\begin{itemize}
\item To build a decentralized voting system where votes are recorded on a public blockchain.
\item To develop a Solidity smart contract that enforces all election rules automatically.
\item To provide a user-friendly Next.js web interface for voters and administrators.
\item To integrate MetaMask wallet authentication for secure, password-free identity verification.
\item To support decentralized image/data storage using IPFS via Pinata.
\item To deploy the system on the Polygon blockchain for low-cost and fast transactions.
\item To demonstrate real-time vote counts and tamper-proof winner determination.
\end{itemize}

\section{Scope of the Project}

The system covers the following key functional areas:

\begin{itemize}
\item Creating and managing elections with configurable time windows.
\item Registering candidates (with name, age, image, and wallet address).
\item Whitelisting voters to participate in a specific election.
\item Allowing registered voters to cast one vote per election.
\item Automatically determining the winner when the election period ends.
\item Providing an admin dashboard for monitoring the election status.
\end{itemize}

The current version is scoped as a testnet deployment and is suitable for college,
organisational, or community-level elections.

\section{Methodology}

The system was developed following the \textbf{Agile Iterative Model}
with the following phases:

\begin{itemize}
\item \textbf{Phase 1 -- Requirement Analysis:} Gathering functional and
      non-functional requirements.
\item \textbf{Phase 2 -- Smart Contract Design:} Writing and testing the
      Solidity \texttt{VotingContract}.
\item \textbf{Phase 3 -- Frontend Development:} Building the Next.js web application
      and integrating Web3 connectivity.
\item \textbf{Phase 4 -- Integration \& Testing:} Connecting all components and
      running unit and integration tests using Hardhat and Chai.
\item \textbf{Phase 5 -- Deployment:} Deploying the smart contract to the Polygon
      Mumbai testnet and hosting the frontend.
\end{itemize}

%=============================================================
%  CHAPTER 2: LITERATURE REVIEW
%=============================================================
\chapter{Literature Review}

\section{Existing Systems}

Several electronic voting and blockchain-based voting systems exist in the literature:

\begin{itemize}
\item \textbf{Voatz (USA):} A mobile blockchain voting app used in West Virginia
      primaries. However, it relies on a private blockchain, limiting public auditability.
\item \textbf{Agora (Sierra Leone):} Used in the 2018 general election. Votes were
      recorded on a blockchain, but the system was not open-source.
\item \textbf{Follow My Vote:} An open-source blockchain voting platform using
      BitShares. Complex to use and not widely deployed.
\item \textbf{Ethereum-based Academic Prototypes:} Multiple research papers propose
      Ethereum smart contract voting systems but lack production-ready frontends.
\end{itemize}

\section{Limitations of Existing Systems}

\begin{itemize}
\item Use of private or permissioned blockchains reduces transparency.
\item Many systems require complex setup and prior blockchain knowledge from voters.
\item Lack of IPFS integration means candidate data is stored in centralised servers.
\item No built-in voter whitelisting mechanism in most prototypes.
\item High transaction fees on Ethereum mainnet make them impractical for common use.
\end{itemize}

\section{Proposed Solution Advantages}

Our Decentralized E-Voting System addresses the above limitations by offering:

\begin{itemize}
\item A \textbf{public blockchain} (Polygon) for full auditability by anyone.
\item A \textbf{voter-friendly Next.js web interface} accessible via MetaMask.
\item \textbf{IPFS-based decentralised storage} for all candidate and voter data.
\item \textbf{Solidity-enforced whitelist} ensuring only authorised voters participate.
\item \textbf{Polygon Layer-2} for low gas costs making it economically viable.
\end{itemize}

%=============================================================
%  CHAPTER 3: SOFTWARE REQUIREMENTS SPECIFICATION (SRS)
%=============================================================
\chapter{Software Requirements Specification}

\section{Functional Requirements}

\begin{enumerate}

\item \textbf{FR1 -- Election Management:}
The system administrator shall be able to create an election by specifying
the election title, organisation name, start time, and end time.

\item \textbf{FR2 -- Candidate Registration:}
The administrator shall register candidates by providing their name, age,
wallet address, and profile image.

\item \textbf{FR3 -- Voter Registration:}
The administrator shall whitelist voters by registering their wallet address,
name, and profile image.

\item \textbf{FR4 -- Vote Casting:}
A whitelisted voter shall be able to cast exactly one vote for any registered candidate
during the active voting window.

\item \textbf{FR5 -- Vote Validation:}
The system shall reject votes from:
\begin{itemize}
  \item Non-whitelisted addresses.
  \item Voters who have already voted.
  \item Addresses outside the active voting window.
\end{itemize}

\item \textbf{FR6 -- Winner Determination:}
The smart contract shall automatically identify and return the candidate with the
highest vote count once the election period ends.

\item \textbf{FR7 -- Live Vote Count:}
Vote counts per candidate shall be publicly queryable from the blockchain at any time.

\item \textbf{FR8 -- Image Upload:}
Candidate and voter images shall be uploaded to IPFS via the Pinata API and
their hashes stored in the smart contract.

\end{enumerate}

\section{Non-Functional Requirements}

\begin{itemize}
\item \textbf{NFR1 -- Security:} All transactions are signed with the user's private key
      via MetaMask. The contract enforces access control using the \texttt{onlyOwner}
      and \texttt{duringVoting} modifiers.
\item \textbf{NFR2 -- Immutability:} Once a vote is cast, it cannot be modified or
      deleted as it is permanently recorded on the blockchain.
\item \textbf{NFR3 -- Transparency:} The smart contract source code and all transactions
      are publicly visible on the Polygon block explorer (PolygonScan).
\item \textbf{NFR4 -- Performance:} Polygon processes blocks in approximately 2 seconds,
      ensuring near real-time vote confirmation.
\item \textbf{NFR5 -- Availability:} The blockchain network operates 24/7 with no
      single point of downtime.
\item \textbf{NFR6 -- Usability:} The web interface uses standard React/Next.js
      components and is accessible on all modern browsers.
\end{itemize}

\section{System Constraints}

\begin{itemize}
\item Voters must have MetaMask installed and a Polygon-compatible wallet.
\item Voters must hold a minimal amount of MATIC tokens to pay transaction gas fees.
\item The election can only be created by the contract owner (admin wallet).
\item Once voting starts, candidate and voter lists are finalized and cannot be altered.
\end{itemize}

%=============================================================
%  CHAPTER 4: SYSTEM DESIGN
%=============================================================
\chapter{System Design}

\section{System Architecture}

The application follows a three-layer decentralised architecture:

\begin{enumerate}
\item \textbf{Presentation Layer:} Next.js web application running in the user's browser.
\item \textbf{Business Logic Layer:} Solidity Smart Contract deployed on the Polygon blockchain.
\item \textbf{Storage Layer:} Polygon blockchain (votes and state) + IPFS via Pinata (media).
\end{enumerate}

\imgbox{13}{8}{System Architecture Diagram -- three-tier: Browser, Smart Contract, Blockchain + IPFS}

\section{Smart Contract Design}

The core \texttt{VotingContract.sol} contains the following \textbf{data structures}:

\begin{longtable}{|l|l|l|}
\hline
\textbf{Struct} & \textbf{Field} & \textbf{Type} \\
\hline
\endhead
Candidate & candidateId & uint256 \\
Candidate & name       & string \\
Candidate & age        & string \\
Candidate & image      & string \\
Candidate & voteCount  & uint256 \\
Candidate & \_address  & address \\
Candidate & ipfs       & string \\
\hline
Voter & voterId       & uint256 \\
Voter & name          & string \\
Voter & image         & string \\
Voter & voterAddress  & address \\
Voter & allowed       & uint256 \\
Voter & voted         & bool \\
Voter & vote          & uint256 \\
Voter & ipfs          & string \\
\hline
Election & electionTitle      & string \\
Election & organizationTitle  & string \\
Election & startTime          & uint256 \\
Election & endTime            & uint256 \\
\hline
\end{longtable}

\section{Key Smart Contract Functions}

\begin{longtable}{|l|l|l|}
\hline
\textbf{Function} & \textbf{Access} & \textbf{Description} \\
\hline
\endhead
createElection()      & Owner Only  & Sets title, org, start \& end time \\
registerCandidate()   & Owner Only  & Adds a candidate to the election \\
giveVoterRight()      & Owner Only  & Whitelists a voter address \\
vote()                & Voter Only  & Casts one vote for a candidate \\
getWinner()           & Public/View & Returns candidate with highest votes \\
isVotingActive()      & Public/View & Returns true if election is live \\
getRemainingTime()    & Public/View & Returns seconds until election ends \\
getAllCandidates()    & Public/View & Returns all registered candidates \\
getAllVoters()        & Public/View & Returns all registered voters \\
\hline
\end{longtable}

\section{Smart Contract Modifiers}

\begin{longtable}{|l|l|}
\hline
\textbf{Modifier} & \textbf{Enforces} \\
\hline
\endhead
onlyOwner   & Caller must be the contract deployer \\
electionExists & An election must have been created \\
duringVoting & Current time must be within voting window \\
\hline
\end{longtable}

\section{Database Design (On-chain State)}

\begin{longtable}{|l|l|l|}
\hline
\textbf{Mapping / Array} & \textbf{Key} & \textbf{Value} \\
\hline
\endhead
candidates         & address & Candidate struct \\
voters             & address & Voter struct \\
candidateAddresses & index   & address \\
voterAddresses     & index   & address \\
\hline
\end{longtable}

\section{Use Case Diagram}

\imgbox{13}{9}{Use Case Diagram -- Admin: creates election, adds candidates, adds voters; Voter: connects wallet, casts vote; Public: views results}

\section{Sequence Diagram -- Voting Flow}

\imgbox{13}{9}{Sequence Diagram -- Voter connects MetaMask, calls vote() on contract, blockchain records transaction, confirmation returned to frontend}

\section{Frontend Page Design}

\begin{table}[H]
\centering
\begin{tabular}{|l|l|}
\hline
\textbf{Page} & \textbf{Description} \\
\hline
Home (index.js)          & Displays election details and candidate listing \\
Candidate Page           & Admin interface to register new candidates \\
Voter Page               & Admin interface to whitelist new voters \\
Vote Page                & Active voter interface to cast a ballot \\
Results Page             & Displays live vote counts and winner \\
\hline
\end{tabular}
\caption{Frontend Page Overview}
\end{table}

%=============================================================
%  CHAPTER 5: IMPLEMENTATION
%=============================================================
\chapter{Implementation}

\section{Hardware Requirements}

\begin{table}[H]
\centering
\begin{tabular}{ll}
\toprule
\textbf{Component} & \textbf{Minimum Specification} \\
\midrule
Processor & Intel i3 / AMD Ryzen 3 or higher \\
RAM       & 4 GB minimum (8 GB recommended) \\
Storage   & 2 GB free disk space \\
Network   & Broadband internet connection \\
Browser   & Google Chrome (with MetaMask extension) \\
\bottomrule
\end{tabular}
\caption{Hardware Requirements}
\end{table}

\section{Software Tools and Technologies}

\begin{table}[H]
\centering
\begin{tabular}{lll}
\toprule
\textbf{Software / Tool} & \textbf{Version} & \textbf{Purpose} \\
\midrule
Node.js      & v18.12.1  & JavaScript runtime environment \\
Next.js      & v12.2.1   & Frontend React framework \\
Solidity     & \^{}0.8.9 & Smart contract programming language \\
Hardhat      & \^{}2.10.0 & Smart contract dev, testing, deployment \\
Ethers.js    & \^{}5.6.9 & Blockchain interaction library \\
MetaMask     & Latest    & Browser wallet \& transaction signer \\
Web3Modal    & \^{}1.9.8 & Wallet connection modal \\
Pinata       & API v1    & IPFS pinning service for file upload \\
Polygon Mumbai & Testnet & Blockchain network for deployment \\
VS Code      & Latest    & Integrated Development Environment \\
\bottomrule
\end{tabular}
\caption{Software Tools}
\end{table}

\section{Smart Contract Implementation}

The \texttt{VotingContract.sol} is the core of the application.
Key implementation decisions include:

\begin{itemize}
\item The \texttt{owner} is set to the deployer's address in the constructor,
      and all administrative functions are protected by the \texttt{onlyOwner} modifier.
\item Candidate and voter data are stored in Solidity \texttt{mapping(address => Struct)}
      for O(1) lookup efficiency.
\item Time enforcement uses \texttt{block.timestamp} compared against the stored
      \texttt{startTime} and \texttt{endTime} values.
\item The \texttt{voted} boolean in the \texttt{Voter} struct provides the
      double-vote prevention mechanism.
\item Events (\texttt{ElectionCreated}, \texttt{CandidateCreated}, \texttt{VoteCasted})
      are emitted for each state change, enabling frontend real-time updates.
\end{itemize}

\imgbox{13}{7}{Screenshot: VotingContract.sol code in VS Code or Remix IDE}

\section{Frontend Implementation}

\subsection{Wallet Connection}

MetaMask wallet connection is handled via \texttt{Web3Modal} and \texttt{Ethers.js}.
The wallet address is stored in the React Context API (\texttt{VoterContext}) and
shared across all pages without unnecessary re-renders.

\imgbox{13}{7}{Screenshot: MetaMask wallet connection modal on the website}

\subsection{Voting Interface}

The vote page displays all registered candidates fetched live from the blockchain.
When a voter clicks ``Vote'', \texttt{Ethers.js} sends a transaction to the
\texttt{vote()} function with the candidate's address as the argument.
MetaMask prompts the user to sign and confirm the transaction.

\imgbox{13}{7}{Screenshot: Vote casting page showing candidate cards}

\subsection{IPFS Image Upload}

When adding a candidate or voter, the admin selects a profile image.
The image is uploaded to Pinata's IPFS pinning service via their REST API.
The returned IPFS CID (content hash) is stored in the smart contract field \texttt{ipfs}.
Images are retrieved using the Pinata IPFS gateway URL.

\imgbox{13}{6}{Screenshot: Candidate registration page with image upload}

\section{Deployment}

\begin{table}[H]
\centering
\begin{tabular}{ll}
\toprule
\textbf{Item} & \textbf{Details} \\
\midrule
Network         & Polygon Mumbai Testnet \\
Deployment Tool & Hardhat deploy script \\
Config File     & \texttt{hardhat.config.js} \\
ABI Location    & \texttt{artifacts/contracts/VotingContract.json} \\
Contract Address & [Insert deployed contract address here] \\
Block Explorer  & \url{https://mumbai.polygonscan.com} \\
\bottomrule
\end{tabular}
\caption{Deployment Details}
\end{table}

\imgbox{13}{6}{Screenshot: Hardhat deployment terminal output showing contract address}

%=============================================================
%  CHAPTER 6: TESTING
%=============================================================
\chapter{Testing}

\section{Testing Strategy}

The project was tested at multiple levels:

\begin{itemize}
\item \textbf{Unit Testing:} Individual smart contract functions were tested in
      isolation using the Hardhat + Chai testing framework.
\item \textbf{Integration Testing:} The frontend was connected to the locally
      deployed contract and tested end-to-end.
\item \textbf{System Testing:} Full deployment on Polygon Mumbai testnet was
      validated using real MetaMask transactions.
\item \textbf{User Acceptance Testing (UAT):} Peer testing was performed to
      validate the UI and overall voter experience.
\end{itemize}

\section{Unit Test Cases}

\begin{longtable}{|l|l|l|l|}
\hline
\textbf{Test ID} & \textbf{Test Description} & \textbf{Expected Result} & \textbf{Status} \\
\hline
\endhead
TC01 & Deploy contract & Owner set to deployer address & Pass \\
TC02 & Create election & Election struct stored correctly & Pass \\
TC03 & Register candidate & Candidate added to mapping & Pass \\
TC04 & Whitelist voter & Voter \texttt{allowed} set to 1 & Pass \\
TC05 & Cast vote during window & Vote count incremented & Pass \\
TC06 & Vote before start time & Transaction reverted  & Pass \\
TC07 & Vote after end time & Transaction reverted & Pass \\
TC08 & Non-whitelisted vote & Transaction reverted & Pass \\
TC09 & Double vote attempt & Transaction reverted & Pass \\
TC10 & Get winner after end & Correct candidate returned & Pass \\
TC11 & Get remaining time & Returns correct seconds & Pass \\
TC12 & isVotingActive during window & Returns true & Pass \\
TC13 & Admin registers duplicate candidate & Transaction reverted & Pass \\
TC14 & Non-owner calls registerCandidate & Transaction reverted & Pass \\
\hline
\end{longtable}

\imgbox{13}{6}{Screenshot: Hardhat test results terminal showing all tests passing}

\section{Frontend Test Cases}

\begin{table}[H]
\centering
\begin{tabular}{|l|l|l|l|}
\hline
\textbf{Test ID} & \textbf{Action} & \textbf{Expected Result} & \textbf{Status} \\
\hline
FT01 & Connect MetaMask & Wallet address displayed in navbar & Pass \\
FT02 & Admin adds candidate & Candidate appears on home page & Pass \\
FT03 & Admin adds voter & Voter appears in voter list & Pass \\
FT04 & Voter casts vote & MetaMask popup shown, vote confirmed & Pass \\
FT05 & Non-voter tries to vote & Error message shown & Pass \\
FT06 & View live vote counts & Counts update in real time & Pass \\
FT07 & Upload image to IPFS & Image displayed using IPFS URL & Pass \\
\hline
\end{tabular}
\caption{Frontend Test Cases}
\end{table}

%=============================================================
%  CHAPTER 7: RESULTS AND DISCUSSION
%=============================================================
\chapter{Results and Discussion}

\section{System Overview}

The Decentralized E-Voting System was successfully implemented and deployed
on the Polygon Mumbai testnet. The system demonstrates all core requirements
including secure voting, automatic result computation, and tamper-proof data storage.

\section{Application Screenshots}

\subsection{Home Page -- Election Dashboard}

The home page displays the current election name, organization, and countdown
timer to the election end. All registered candidates are listed with their names,
ages, and current vote counts.

\imgbox{13}{8}{Screenshot: Home/Dashboard page of the E-Voting DApp}

\subsection{Candidate Registration (Admin)}

The admin panel allows the election administrator to add new candidates by
entering their wallet address, name, age, and uploading a profile image.
The image is sent to IPFS and the hash stored on-chain.

\imgbox{13}{8}{Screenshot: Candidate registration form -- admin view}

\subsection{Voter Registration (Admin)}

Similar to candidate registration, the admin can whitelist voters by their
wallet address. Once registered, a voter can log in using MetaMask and cast a vote.

\imgbox{13}{8}{Screenshot: Voter registration form -- admin view}

\subsection{Vote Casting Interface}

Registered voters see the list of candidates. Clicking ``Vote'' triggers a
MetaMask transaction. After confirmation, the voter's status updates to
\texttt{voted = true} on-chain and the candidate's vote count increments.

\imgbox{13}{8}{Screenshot: Voter casting ballot -- MetaMask confirmation popup}

\subsection{Election Results}

After the election window closes, the smart contract's \texttt{getWinner()}
function returns the winning candidate. The results page displays the final
vote count for all candidates and highlights the winner.

\imgbox{13}{8}{Screenshot: Election results page showing winner and final vote counts}

\section{Performance Observations}

\begin{table}[H]
\centering
\begin{tabular}{lll}
\toprule
\textbf{Operation} & \textbf{Network} & \textbf{Time / Gas Cost} \\
\midrule
Deploy Contract   & Polygon Mumbai & Approx. 2--5 seconds / \textasciitilde{}0.01 MATIC \\
Register Candidate & Polygon Mumbai & Approx. 2 seconds / \textasciitilde{}0.001 MATIC \\
Cast Vote          & Polygon Mumbai & Approx. 2 seconds / \textasciitilde{}0.001 MATIC \\
View Results (read) & Any            & Instant / Free \\
\bottomrule
\end{tabular}
\caption{System Performance}
\end{table}

%=============================================================
%  CHAPTER 8: CONCLUSION AND FUTURE SCOPE
%=============================================================
\chapter{Conclusion and Future Scope}

\section{Conclusion}

The \textbf{Decentralized E-Voting System} was successfully designed, developed, tested,
and deployed as a full-stack blockchain application. The project demonstrates that
it is technically feasible to replace centralised voting infrastructure with a
transparent, trustless, and tamper-proof system built on public blockchain technology.

The key achievements of this project are:

\begin{itemize}
\item A production-ready Solidity smart contract with comprehensive access control,
      voter whitelisting, and automatic result computation.
\item A fully functional Next.js web frontend connected to the blockchain via
      Ethers.js and MetaMask.
\item Cost-effective deployment on the Polygon Layer-2 network with gas costs
      reduced by over 99\% compared to Ethereum mainnet.
\item Decentralized media storage using IPFS through the Pinata pinning service.
\item Comprehensive smart contract testing with 14 unit test cases all passing.
\end{itemize}

This project proves that blockchain-based e-voting is not merely a theoretical
concept but a working, deployable solution that can serve organisations, colleges,
or local government bodies.

\section{Limitations}

\begin{itemize}
\item Voters require a MetaMask wallet and MATIC tokens, creating a technical barrier
      for non-technical users.
\item The current version does not support voter identity verification against
      government records.
\item Voter anonymity is limited since wallet addresses are publicly visible on-chain.
\end{itemize}

\section{Future Enhancements}

\begin{itemize}
\item \textbf{Zero-Knowledge Proofs (ZK-SNARKs):} Implement cryptographic
      voter anonymity using zk-SNARK circuits -- prove registration without revealing identity.
\item \textbf{Aadhaar / Government ID Integration:} Automate voter eligibility
      verification using national identity APIs to prevent fake registrations.
\item \textbf{Mobile Application:} Develop a dedicated Android and iOS app using
      React Native for wider accessibility.
\item \textbf{DAO Governance:} Replace the single admin model with a
      Decentralized Autonomous Organization (DAO) structure for multi-authority elections.
\item \textbf{IPFS for Ballot Storage:} Store the complete encrypted ballot on IPFS
      for an additional layer of auditability.
\item \textbf{Multi-Election Support:} Extend the contract to support multiple
      simultaneous elections within one deployment.
\end{itemize}

%=============================================================
%  REFERENCES
%=============================================================
\chapter*{References}
\addcontentsline{toc}{chapter}{References}

\begin{enumerate}
\item Solidity Documentation. \textit{Solidity Language Reference}. \url{https://docs.soliditylang.org}

\item Hardhat. \textit{Hardhat Development Environment Documentation}. \url{https://hardhat.org/docs}

\item Ethers.js. \textit{Ethers.js v5 Documentation}. \url{https://docs.ethers.org/v5/}

\item Polygon. \textit{Polygon Developer Documentation}. \url{https://wiki.polygon.technology}

\item Pinata. \textit{Pinata IPFS Pinning API Reference}. \url{https://docs.pinata.cloud}

\item Next.js. \textit{Next.js Framework Documentation}. \url{https://nextjs.org/docs}

\item Web3Modal. \textit{Web3Modal Documentation}. \url{https://github.com/Web3Modal/web3modal}

\item IPFS. \textit{InterPlanetary File System Documentation}. \url{https://docs.ipfs.tech}

\item Nakamoto, S. (2008). \textit{Bitcoin: A Peer-to-Peer Electronic Cash System}.

\item Buterin, V. (2014). \textit{Ethereum: A Next-Generation Smart Contract and
      Decentralised Application Platform}. Ethereum White Paper.

\item Kshetri, N. and Voas, J. (2018). Blockchain-Enabled E-Voting.
      \textit{IEEE Software}, 35(4), 95--99.

\item OpenZeppelin. \textit{OpenZeppelin Smart Contract Standards}.
      \url{https://docs.openzeppelin.com}
\end{enumerate}

%=============================================================
%  APPENDIX
%=============================================================
\chapter*{Appendix}
\addcontentsline{toc}{chapter}{Appendix}

\section*{Appendix A: Smart Contract Source Code}

The complete source code of \texttt{VotingContract.sol} is available at:\\
\url{https://github.com/[your-github]/e-voting-dapp}

\imgbox{13}{10}{Screenshot: Full VotingContract.sol source code}

\section*{Appendix B: Hardhat Configuration}

\texttt{hardhat.config.js} configures the compiler version and Polygon Mumbai
network RPC endpoint using environment variables from the \texttt{.env} file.

\imgbox{13}{6}{Screenshot: hardhat.config.js file}

\section*{Appendix C: Environment Variables}

Required environment variables stored in \texttt{.env}:

\begin{table}[H]
\centering
\begin{tabular}{ll}
\toprule
\textbf{Variable} & \textbf{Purpose} \\
\midrule
NEXT\_PUBLIC\_CONTRACT\_ADDRESS & Deployed smart contract address \\
PINATA\_API\_KEY               & Pinata service API key \\
PINATA\_SECRET\_API\_KEY       & Pinata service secret \\
POLYGON\_MUMBAI\_RPC\_URL      & Polygon Mumbai RPC endpoint \\
PRIVATE\_KEY                  & Admin wallet private key for deployment \\
\bottomrule
\end{tabular}
\caption{Environment Variables}
\end{table}

\section*{Appendix D: Sample Deployed Transaction}

\imgbox{13}{6}{Screenshot: PolygonScan transaction record of a vote being cast}

\end{document}
